% =====================================================
% PART 2: DATA PREPROCESSING
% =====================================================
\chapter{Data Preprocessing}

The preprocessing pipeline consists of 6 steps, each justified by the findings of Part 1.

\section{Excel error correction and categorical variable handling}

\begin{lstlisting}[language=Python]
# Replace Excel error string with NaN
df = df.replace('#NOM?', np.nan)

# Categorical variables: NaN -> 'UNKNOWN'
for col in ['Protocole', 'Service_cible']:
    df[col] = df[col].fillna('UNKNOWN')
\end{lstlisting}

The category \texttt{UNKNOWN} avoids removing rows (up to 8\% of the dataset) and treats missing information as a separate category.

\section{Conversion and imputation of numerical variables}

\begin{lstlisting}[language=Python]
colonnes_exclues = ['IP_source', 'IP_destination', 'Protocole', 'Service_cible',
                     'Pays_source', 'Heure_connexion', 'Etiquete']
colonnes_quantitatives = [c for c in df.columns if c not in colonnes_exclues]

for col in colonnes_quantitatives:
    df[col] = pd.to_numeric(df[col], errors='coerce')
    moyenne = df[col].mean()
    df[col] = df[col].fillna(moyenne if pd.notna(moyenne) else 0)
\end{lstlisting}

\textbf{Justification:} mean imputation preserves the 103,000 observations with limited bias due to the low percentage of missing values. A fallback to 0 is used if a column is entirely empty.

\section{Removal of irrelevant columns}

\begin{lstlisting}[language=Python]
df = df.drop(columns=['Pays_source', 'Heure_connexion'], errors='ignore')
\end{lstlisting}

\texttt{Pays\_source} may introduce unwanted geographic bias. \texttt{Heure\_connexion} requires temporal modeling outside the scope of this work.

\section{Encoding categorical variables}

\begin{lstlisting}[language=Python]
from sklearn.preprocessing import LabelEncoder

encodeur_proto = LabelEncoder()
encodeur_service = LabelEncoder()
df['Protocole'] = encodeur_proto.fit_transform(df['Protocole'].astype(str))
df['Service_cible'] = encodeur_service.fit_transform(df['Service_cible'].astype(str))
\end{lstlisting}

\textbf{Label Encoding} is preferred over One-Hot Encoding to avoid feature explosion. Random Forest handles encoded categorical variables efficiently without requiring feature scaling.

\section{X / y separation}
\label{sec:separation}

\begin{lstlisting}[language=Python]
from sklearn.model_selection import train_test_split

if '0' in df.columns:
    df = df.drop(columns=['0'])

y = df['Etiquete']
X = df.drop(columns=['Etiquete', 'IP_source', 'IP_destination'], errors='ignore')

print(X.shape, y.shape)
print(list(X.columns))
\end{lstlisting}

\textbf{IP removal:} IP addresses are identifiers, not behavioral features. Keeping them could lead to data leakage, where the model memorizes specific IPs instead of learning general patterns. They are still kept in \texttt{df} for reporting purposes (Part 4).

\textbf{Result:} $X$ contains \textbf{17 explanatory variables} for 103,000 samples; $y$ is the target vector.

\section{Normalization: not applied}

\textbf{Random Forest is invariant to monotonic transformations}, as decision trees rely on threshold splits. Therefore, scaling methods such as \texttt{StandardScaler} or \texttt{MinMaxScaler} are unnecessary.

\section{Train/test split}

\begin{lstlisting}[language=Python]
X_train, X_test, y_train, y_test = train_test_split(
    X, y, test_size=0.30, random_state=42
)
print(X_train.shape[0], X_test.shape[0])  # 72100 / 30900
\end{lstlisting}

A 70/30 split is used with \texttt{random\_state=42} for reproducibility. The dataset is now ready for model training (Part 3).